\documentclass[11pt,a4paper]{article}
\usepackage[utf8]{inputenc}
\usepackage[T1]{fontenc}
\usepackage[brazil]{babel}
\usepackage[top=1.8cm, bottom=1.8cm, left=2.0cm, right=2.0cm]{geometry}
\usepackage{lmodern}
\usepackage{array}
\usepackage{tabularx}
\usepackage{booktabs}
\usepackage{colortbl}
\usepackage{xcolor}
\usepackage{enumitem}
\usepackage{amsmath}
\usepackage{titlesec}
\usepackage{fancyhdr}

\definecolor{uefsgray}{HTML}{EDEDED}
\definecolor{uefslight}{HTML}{F7F7F7}
\definecolor{uefsdark}{HTML}{1A1A1A}

\titleformat{\section}{\normalfont\bfseries\large}{}{0em}{\MakeUppercase}[\titlerule]
\titlespacing*{\section}{0pt}{12pt}{6pt}

\pagestyle{fancy}
\fancyhf{}
\renewcommand{\headrulewidth}{0pt}
\fancyfoot[L]{\footnotesize UEFS $\cdot$ DCHF $\cdot$ CHF662 Análise Ambiental}
\fancyfoot[R]{\footnotesize Página \thepage}

\setlength{\parindent}{0pt}
\setlength{\parskip}{6pt}

\begin{document}

\begin{center}
  {\bfseries\large UNIVERSIDADE ESTADUAL DE FEIRA DE SANTANA}\\[2pt]
  {\bfseries\normalsize DEPARTAMENTO DE CIÊNCIAS HUMANAS E FILOSOFIA}\\[2pt]
  {\bfseries\normalsize CURSO DE GEOGRAFIA}\\[8pt]
  \rule{\textwidth}{1.2pt}\\[6pt]
  {\bfseries\large ATIVIDADE DE ESTUDO (AULA 12.1)}
\end{center}

\vspace{2pt}

\begin{center}
\renewcommand{\arraystretch}{1.25}
\begin{tabularx}{\textwidth}{|>{\bfseries}p{3.8cm}|X|>{\bfseries}p{2.9cm}|X|}
\hline
\rowcolor{uefsgray}
Componente Curricular & CHF662 Análise Ambiental (60h/a) & Período Letivo & 2026/2 \\
\hline
\rowcolor{uefsgray}
Docente & Prof. Dr. Luiz Diego Vidal Santos & Referência & Aula 12.1 e Capítulo 11 \\
\hline
\end{tabularx}
\end{center}

\vspace{2pt}

\begin{tabularx}{\textwidth}{|X|p{3.2cm}|p{3.5cm}|}
\hline
\textbf{Discente:} & \textbf{Turma:} & \textbf{Data:} \quad / \quad / \quad \\
\hline
\end{tabularx}

\vspace{6pt}

\section{Onde estudar}

\noindent\fcolorbox{uefsdark}{uefslight}{%
\begin{minipage}{\dimexpr\textwidth-2\fboxsep-2\fboxrule\relax}
\small
Baseie suas respostas nos slides da \textbf{Aula 12.1 (Qualidade do Solo e Capacidade de Uso das Terras)} e no \textbf{Capítulo 11} da apostila (páginas \textbf{264 a 283}). Consulte as seções sobre conceito de qualidade do solo e tríade do MDS (p.\,265), propriedades inerentes e dinâmicas, indicadores físicos, químicos e biológicos do MDS (p.\,266--272) e amostragem com área de referência. Atividade sem cálculos numéricos; conclua antes de consultar o gabarito.
\end{minipage}%
}

\section{Questões}

\begin{enumerate}[leftmargin=*, itemsep=6pt]
  \item Explique o conceito contemporâneo de qualidade do solo e os motivos pelos quais a avaliação da funcionalidade edáfica excede a análise exclusiva de fertilidade química.

  \item Diferencie propriedades inerentes e propriedades dinâmicas do solo, apresentando dois exemplos de cada grupo e indicando suas escalas temporais de alteração.

  \item Descreva quatro funções ecossistêmicas desempenhadas pelo solo na paisagem, relacionando-as à regulação hídrica, à ciclagem biogeoquímica e à estabilidade física.

  \item Distinga indicadores físicos, químicos e biológicos da qualidade do solo, apresentando um exemplo representativo de cada categoria e o processo edáfico que ele expressa.

  \item Explique por que a comparação com uma área de referência (como um fragmento de vegetação nativa na mesma classe pedológica, textura e posição geomorfológica) condiciona a validade do diagnóstico de degradação.

  \item Defina o Conjunto Mínimo de Dados (MDS) e indique quais critérios operacionais e ecológicos orientam a seleção de seus indicadores.

  \item Assinale a alternativa correta sobre a avaliação da qualidade do solo,
  \begin{enumerate}[label=(\Alph*), leftmargin=1.5em, itemsep=2pt, topsep=3pt]
    \item Deve considerar somente a fertilidade e o pH.
    \item Deve integrar propriedades físicas, químicas e biológicas, considerando o uso e o contexto ambiental.
    \item Deve utilizar os mesmos valores de referência para todos os tipos de solo.
    \item Deve dispensar observações de campo quando houver análise laboratorial.
    \item Deve avaliar apenas a produtividade agrícola.
  \end{enumerate}

  \item Uma área submetida a tráfego intenso apresenta solo compactado, baixa cobertura vegetal e reduzida taxa de infiltração. Explique o que esses sinais indicam sobre a arquitetura porosa do solo e quais processos hidrossedimentológicos tendem a intensificar-se durante eventos pluviométricos.

  \item Explique por que o carbono da biomassa microbiana ($C\text{-BM}$) e a respiração basal do solo ($RBS$) atuam como indicadores precoces de alterações na qualidade edáfica em comparação ao carbono orgânico total.

  \item Para um diagnóstico de recuperação ambiental, selecione quatro indicadores de qualidade do solo distribuídos entre as dimensões física, química e biológica, justificando a escolha com base nas funções ecossistêmicas monitoradas.
\end{enumerate}

\end{document}
